\subsection{LTOF左右対応の確認と片側ヒットの診断}
先の隣接ch ANDとFlat ORの比較を、同一LTOFセグメントの左右ANDとの比較と説明していた点を訂正する。marq04の保存設定のAND組は(16,17)、(18,19)、(20,21)、(22,23)、(24,25)、(27,28)であり、それぞれ異なるセグメントの左同士または右同士であった。性能値と入力はその保存条件での実測値として保持するが、左右読み落としの測定には使わない。

Artemisの2024年用tof.mapと、今回のchannel-map用CSV（20260311版を現行API形式へ変換）およびmarq04の20260909版は、LTOF全12読み出しのFEM、セグメント、左右、チャンネルが一致した。ArtemisのセグメントIDは0--5であり、以下の図表では1--6と表記する。正しい左右組は(16,22)、(19,25)、(17,23)、(20,27)、(18,24)、(21,28)である。マップ不一致ではなく、LogicFilterに直接指定された組合せの問題であった。

marq06のトラッキング用launcher（更新日時2026年9月11日）はmq-param/triggerlogic.sh（9月10日）を呼び、Flat ORの式が有効であった。隣接ch AND版はコメントアウトされていた。いずれのサーバーの保存設定も変更していない。

正しい同一セグメント左右AND後ORを新たな条件として、marq07で結合run000400とrun000668の先頭20,000 TFを確認した。HLFは既存TOFを通したが、ここでは受信前のSlice数のみを数え、性能測定としては扱わない。幅30、LTOF補正$-12$、Slice窓$[-250,250]$は従来と同じ4 ns単位である。全プロセス終了0と、LogicFilterの正のTotalHits列とTFごとの正のSlice数列の一致を確認した。
\begin{table}[htbp]\centering\small
\caption{marq07、全20,000 TF、41.94304 sに対するトリガー数。0 Slice TFも含めた平均。}
\begin{tabular}{rlrrr}\hline
run & 条件 & 総トリガー数 & 平均Slice/TF & 平均レート[kHz] \\\hline
400 & marq04隣接ch AND & 156,792 & 7.8396 & 3.738 \\
400 & 同一seg左右AND & 2,183,856 & 109.1928 & 52.067 \\
400 & Flat OR & 2,957,669 & 147.8835 & 70.516 \\
668 & marq04隣接ch AND & 151,862 & 7.5931 & 3.621 \\
668 & 同一seg左右AND & 2,143,307 & 107.1654 & 51.100 \\
668 & Flat OR & 2,919,747 & 145.9873 & 69.612 \\
\hline\end{tabular}\end{table}

生データ診断はLF/HLFで選別する前のデータを直接復号し、marq07の最大124スレッドで行った。UTOF左ヒットごとに左右差10 ns以内の最寄り右ヒットを1つ選び、その平均を参照時刻とした。LTOFの存在を参照条件に要求しない。HBF境界から600 ns以内の参照は除外した。同じヒットが複数参照に対応する場合があるため、以下はヒット・参照の組合せによる頻度であり固有検出効率ではない。

LTOF各側の参照からの時間には約45--50 nsにピークがある。各segで$[-72,168)$ nsと$[20,80)$ nsの2つの窓について、Lのみ、Rのみ、両方、どちらもなしを数えた。左右の時間差分布は窓内の全組合せで作成し、10、30、120 ns以内の左右ペアが存在する参照数も保存した。広い窓と同じ240 ns幅の前後サイドバンド$[-312,-72)$、$[168,408)$ nsも数えた。サイドバンドにも相関を持つLTOFヒットがあり、単純差引きを固有効率と解釈しない。

TOTはV3 leadingワードに格納された値を1024で割ってnsへ換算した。Artemisのds-qtc設定も時間・電荷に共通の約1/1024変換を用いている。TOTは光量そのものを較正した量ではない。片側のみと両側のTOT分布は面積で規格化し、頻度の差は別表・CSVで示す。図は線形軸、200ビン、グリッド付きとした。
\begin{table}[htbp]\centering\small
\caption{run000668、LTOF各側の時刻がUTOF参照から$[20,80)$ nsにある場合。割合の分母は各segで少なくとも片側がある参照数。TOTは対応するヒット・参照組合せの平均[ns]。}
\begin{tabular}{rrrrrrrr}\hline
seg & Lのみ[\%] & Rのみ[\%] & 両側[\%] & LのみTOT & 両側L TOT & RのみTOT & 両側R TOT \\\hline
1 & 49.13 & 8.84 & 42.03 & 9.14 & 31.13 & 13.91 & 27.60 \\
2 & 40.54 & 23.33 & 36.13 & 14.17 & 20.64 & 12.32 & 16.98 \\
3 & 20.39 & 18.63 & 60.98 & 21.23 & 20.80 & 19.07 & 18.91 \\
4 & 10.64 & 11.31 & 78.05 & 23.18 & 23.86 & 23.91 & 24.24 \\
5 & 17.65 & 14.84 & 67.51 & 13.34 & 27.60 & 15.74 & 25.05 \\
6 & 17.64 & 12.86 & 69.50 & 9.44 & 29.45 & 18.77 & 33.72 \\
\hline\end{tabular}\end{table}

正しい左右ANDのレートはFlat ORの約73--74\%であり、隣接ch ANDとの約19倍差を光量不足に帰属する解釈は成立しない。主な照射領域のseg3、4にも片側ヒットはあるが、片側ヒットのTOTが両側時のTOTより一様に小さいわけではない。一方seg1、5、6などでは片側時のTOTが小さい傾向がある。低光量や閾値による読み落としは候補として残るが、位置依存、ゲイン・閾値、別粒子、背景、窓幅の影響と分離できておらず、光量不足が主因とは断定しない。

今回対象のUTOF/LTOF FEMには独立なSPL開始・終了ワードを確認できず、HBFのflagをスピル状態として検証していない。したがって、時系列の低レート区間を確定したspill offとは呼ばない。次の切り分けには独立なスピル信号、通過位置、閾値・ゲイン情報との照合が必要である。

\begin{figure}[htbp]\centering\makebox[\linewidth][c]{\includegraphics[pagebox=cropbox,width=.9\linewidth,height=.65\textheight,keepaspectratio]{05-hlf/ltof-readout-diagnosis/run668-trigger-time.pdf}}
\caption{3条件のトリガーレート時系列。TF中央時刻への0.1 s割当、末尾0.04304 sは実時間で規格化。}
\label{fig:ltof-diagnosis-run668-trigger-time}
\end{figure}

\begin{figure}[htbp]\centering\makebox[\linewidth][c]{\includegraphics[pagebox=cropbox,width=.9\linewidth,height=.65\textheight,keepaspectratio]{05-hlf/ltof-readout-diagnosis/run668-seg3-time.pdf}}
\caption{run000668、seg3の各側のUTOF参照からの時間。}
\label{fig:ltof-diagnosis-run668-seg3-time}
\end{figure}

\begin{figure}[htbp]\centering\makebox[\linewidth][c]{\includegraphics[pagebox=cropbox,width=.9\linewidth,height=.65\textheight,keepaspectratio]{05-hlf/ltof-readout-diagnosis/run668-seg4-time.pdf}}
\caption{run000668、seg4の各側のUTOF参照からの時間。}
\label{fig:ltof-diagnosis-run668-seg4-time}
\end{figure}

\begin{figure}[htbp]\centering\makebox[\linewidth][c]{\includegraphics[pagebox=cropbox,width=.9\linewidth,height=.65\textheight,keepaspectratio]{05-hlf/ltof-readout-diagnosis/run668-seg3-lr.pdf}}
\caption{run000668、seg3の左右時間差。各窓内の全左右組合せ。}
\label{fig:ltof-diagnosis-run668-seg3-lr}
\end{figure}

\begin{figure}[htbp]\centering\makebox[\linewidth][c]{\includegraphics[pagebox=cropbox,width=.9\linewidth,height=.65\textheight,keepaspectratio]{05-hlf/ltof-readout-diagnosis/run668-seg3-tot-tight.pdf}}
\caption{run000668、seg3の20--80 ns窓におけるTOT分布。各曲線は全ヒット・参照組合せで面積規格化、表示範囲外も規格化に含める。}
\label{fig:ltof-diagnosis-run668-seg3-tot-tight}
\end{figure}

\begin{figure}[htbp]\centering\makebox[\linewidth][c]{\includegraphics[pagebox=cropbox,width=.9\linewidth,height=.65\textheight,keepaspectratio]{05-hlf/ltof-readout-diagnosis/run668-seg4-tot-tight.pdf}}
\caption{run000668、seg4の20--80 ns窓におけるTOT分布。規格化はseg3と同じ。}
\label{fig:ltof-diagnosis-run668-seg4-tot-tight}
\end{figure}
